% =============================================================================
% Cap. 1 — Introdução (intercalar conforme guia Pestana §3)
% Fonte: docs/scope/proposta.tex (contrato com orientador, 22 mar 2026)
%        + docs/scope/requirements.tex (MoSCoW v1.0)
%        + docs/scope/risks.tex
% =============================================================================

\section{Introdução}
\label{cap:introducao}

O ForensiQ é uma aplicação \emph{web} \emph{mobile-first} que digitaliza
e padroniza o registo, a preservação e a transferência de prova
digital pelos agentes que primeiro chegam a uma cena de crime. Este
relatório intercalar apresenta o estado do projecto à data de 6 de
Maio de 2026, em cumprimento do calendário da Unidade Curricular 21184
da Universidade Aberta. O documento está estruturado em três
capítulos, conforme o Guia de Projecto do orientador: introdução
(presente capítulo), desenho (Capítulo~\ref{cap:desenho}) e
implementação no estado actual (Capítulo~\ref{cap:implementacao}).

\subsection{Descrição do projecto}
\label{sec:descricao}

Os agentes da Polícia de Segurança Pública (PSP) que atuam como
\emph{first responders} em ocorrências com componente digital recorrem
hoje a métodos manuais para registar e preservar a prova: papel,
fotografias avulsas no telemóvel pessoal, anotações manuscritas e
etiquetas em saquetas. O processo é vulnerável em várias dimensões.
A cadeia de custódia torna-se frágil: não existe registo verificável
de quem manuseou a prova, em que momento e em que condições.
A prova carece frequentemente de localização exacta da apreensão e de
\emph{timestamp} sincronizado. As práticas variam entre agentes e entre
divisões, o que dificulta auditoria processual posterior. Em tribunal,
fragilidades de cadeia de custódia podem inviabilizar prova relevante,
com consequências para a justiça material~\citep{casey2011,acpo2012}.

O autor identificou o problema a partir da experiência directa
enquanto agente da PSP e perito digital forense. O tema foi formalizado
como auto-proposta na Fase~0, em alinhamento com o
Projeto~\#38~\citep{pestana2025} da licenciatura, com a delimitação
explícita de que o ForensiQ não realiza aquisição de prova ---
apenas a sua \emph{gestão} e \emph{controlo} após apreensão.

A solução é uma plataforma \emph{web} servida pelo \emph{backend}
Django, acedida por \emph{browser} a partir de \emph{smartphone} ou
computador, sem necessidade de instalação de \emph{software}
especializado. Os registos são imutáveis por desenho, com integridade
verificável por terceiros através de \emph{hash} SHA-256, e a cadeia
de custódia é encadeada criptograficamente para que qualquer
adulteração posterior seja detectável.

\subsection{Âmbito e enquadramento}
\label{sec:ambito}

O âmbito do ForensiQ cobre o ciclo completo \emph{registo $\rightarrow$
preservação $\rightarrow$ transferência} da prova após apreensão.
A versão actual é \emph{digital-only}, conforme registado no ADR-0010,
com a estrutura preparada para módulos forenses adicionais (biológico,
químico, físico) em trabalho futuro. Estão fora do âmbito, de forma
explícita: a aquisição forense de imagens de discos e telemóveis
(\emph{carving}, \texttt{dd}, FTK Imager); a análise de \emph{malware};
a peritagem de ADN ou AFIS; a integração com sistemas internos da PSP,
que requereria autorizações não disponíveis em contexto académico; e
o uso de dados reais de processos judiciais, por questões legais e de
confidencialidade.

\subsection{Necessidades identificadas}
\label{sec:necessidades}

A análise do contexto operacional, validada com o orientador na
proposta de Fase~1, identificou três necessidades transversais que
estruturam todo o desenho da solução:

\begin{enumerate}
    \item \textbf{Integridade verificável.} Qualquer alteração posterior
          a um registo de prova tem de ser detectável por um terceiro
          que tenha acesso ao identificador e aos campos públicos.
          Esta necessidade liga-se directamente ao \S5.4 da norma
          ISO/IEC~27037:2012~\citep{iso27037}.
    \item \textbf{Rastreabilidade auditável.} Cada manuseamento da prova
          --- quem, quando, em que estado --- tem de ficar registado de
          forma imutável, com sequência canónica que sobreviva a
          discrepâncias de relógio entre máquinas. Esta necessidade
          materializa o \S5.5 da norma.
    \item \textbf{Utilidade no terreno.} O instrumento tem de funcionar
          num \emph{smartphone}, com uma só mão, em redes móveis
          variáveis, sem instalação prévia de \emph{software}
          especializado nem formação técnica avançada. Esta necessidade
          é uma condição de adopção: uma plataforma robusta mas
          inutilizável no momento da apreensão não resolve o problema.
\end{enumerate}

A estas três necessidades técnicas acrescem requisitos legais e éticos
não-negociáveis: cumprimento do RGPD, em particular do
Artigo~32.\textordfeminine{}~\citep{rgpd2016}, e ausência de exposição
de dados pessoais em ambiente \emph{open-source}.

\subsection{Objectivos}
\label{sec:objectivos}

\paragraph{Objectivo geral.}
Desenvolver uma plataforma \emph{web} \emph{mobile-first} que
digitalize a recolha, preservação e transferência de prova digital
pelos \emph{first responders} da PSP, com integridade verificável
conforme a norma ISO/IEC~27037:2012, e que seja demonstrável num
\emph{smartphone} em condições reais.

\paragraph{Objectivos específicos.}
Os seis objectivos específicos da proposta correspondem às
funcionalidades F1--F6 do MVP, com critérios de aceitação observáveis.
Cada objectivo é mapeado, no Capítulo~\ref{cap:implementacao}, ao seu
estado actual de implementação e à evidência empírica que o suporta.

\begin{longtable}{@{}L{1.0cm}L{6.0cm}L{7.5cm}@{}}
\toprule
\textbf{Obj.} & \textbf{Funcionalidade} & \textbf{Critério de aceitação observável} \\
\midrule
\endhead
F1 & Autenticação JWT com perfis distintos & Login válido em \(<\)2~s; perfis com permissões verificáveis; erro 401 sem fuga de informação interna. \\
\midrule
F2 & Registo de prova com fotografia e GPS & Criação completa em \(<\)30~s no \emph{smartphone}; identificador único, \emph{timestamp} do servidor, \emph{hash} SHA-256 dos metadados e dos \emph{bytes} da fotografia. \\
\midrule
F3 & Cadeia de custódia com máquina de estados imutável & Transições só na ordem canónica definida; sem retrocesso; \emph{append-only} sem \texttt{UPDATE} nem \texttt{DELETE}; alteração directa na base de dados detectável por verificação de \emph{hash}. \\
\midrule
F4 & Módulo de forense digital (ficha de dispositivo) & Ficha completa associada a uma evidência existente: tipo, marca, modelo, estado, número de série. \\
\midrule
F5 & API REST documentada & Mínimo dez \emph{endpoints} testáveis via Swagger~UI gerado automaticamente. \\
\midrule
F6 & Exportação de relatório de ocorrência em PDF & Documento contendo todos os registos e o histórico de custódia da ocorrência, juntável ao expediente físico. \\
\bottomrule
\end{longtable}

\subsection{Restrições}
\label{sec:restricoes}

A execução do projecto está sujeita a quatro restrições documentadas
em \texttt{docs/scope/risks.tex}: duração de 16~semanas com três
marcos eliminatórios fixos (proposta, intercalar, final); trabalho
estritamente individual sem possibilidade de divisão de tarefas;
dependência de planos gratuitos da Neon (PostgreSQL gerido) e da
Fly.io (alojamento Frankfurt) durante o desenvolvimento, com migração
para \emph{Neon Launch} (\(\sim\)5~USD/mês) antes da defesa pública;
e indisponibilidade de orçamento para serviços comerciais,
\emph{e.g.} \emph{APIs} pagas de OCR ou OSINT.

A estas restrições acrescem condicionantes éticas: o sistema é
desenhado para uso futuro com dados reais, mas é demonstrado
exclusivamente com dados sintéticos gerados por \emph{seed}; o
repositório é público no GitHub e segue o princípio de
``segurança por desenho, transparência por omissão''.

\subsection{Resultados esperados}
\label{sec:resultados}

A verificação do cumprimento dos objectivos é tripla: pelo
funcionamento da aplicação em produção em
\href{https://forensiq.pt}{forensiq.pt}, pelos testes automatizados
no repositório e pela validação externa de postura de segurança.

A aplicação em produção permite a um avaliador externo,
nomeadamente o orientador, executar os fluxos críticos --- criação
de ocorrência, registo de evidência com GPS e fotografia,
transição de custódia, exportação de PDF --- sem necessidade de
configurar ambiente local. Esta capacidade serve de demonstração
assíncrona do MVP e é abordada em
\S\ref{sec:demo}. A \emph{suite} automatizada de testes
unitários e de integração (Capítulo~\ref{cap:implementacao})
exercita os caminhos críticos de imutabilidade, autorização e
encadeamento de \emph{hash}. Os três testes externos de segurança
em \texttt{docs/compliance/external-tests/} (HSTS \emph{Preload}
submetido a Chromium, Mozilla, Edge e Apple; Mozilla HTTP
Observatory~A+; Qualys SSL Labs~A+) validam, com auditoria
independente, a postura criptográfica e os cabeçalhos HTTP.

\subsection{Calendário planeado}
\label{sec:calendario}

A Tabela~\ref{tab:calendario} apresenta o calendário das 16~semanas
do projecto. Os marcos formais (Fase~1 a 25~mar, Fase~2 a 6~mai,
Fase~3 a 24~jun) são fixos. As semanas intermédias foram executadas
com pequenos ajustes face à proposta; em particular, a demo interna
síncrona prevista para a Sem.~7 (28~abr--2~mai) não foi realizada
nessa janela, sendo substituída por demonstração assíncrona em
produção --- assunto desenvolvido em \S\ref{sec:deviations} e
comunicado ao orientador em paralelo a este relatório.

\begin{longtable}{@{}L{1cm}L{2.4cm}L{7.5cm}L{2.6cm}@{}}
\toprule
\textbf{Sem.} & \textbf{Datas} & \textbf{Conteúdo planeado} & \textbf{Marco} \\
\midrule
\endhead
1--2  & 17--28~mar     & Proposta inicial, configuração do repositório. & Proposta (25~mar) \\
\midrule
3--4  & 31~mar--11~abr & MoSCoW final, ADRs 0001--0005, modelo ER, primeiros \emph{models} Django. & \\
\midrule
5--6  & 14--25~abr     & Autenticação JWT em \emph{cookies} HttpOnly (ADR-0009), CSP \emph{nonce-based} (ADR-0007), F2 e F3 funcionais. & \\
\midrule
7     & 28~abr--2~mai  & Auditoria interna de segurança (\texttt{AUDIT\_2026-04-16.md}), correcções \emph{B-C2} e \emph{B-C3}, integração de \emph{triggers} PostgreSQL, ADRs 0006--0010. \emph{Demo interna síncrona não realizada nesta janela.} & Demo interna \\
\midrule
8     & 5--6~mai       & Submissão do relatório intercalar. & \textbf{Intercalar (6~mai)} \\
\midrule
9--10 & 7--16~mai      & Reforço de cobertura de testes; revisão de UX em condições de terreno; eventual demo síncrona com o orientador. & \\
\midrule
11--12 & 19--30~mai    & Polimento da \emph{interface}, módulo extra opcional (\emph{e.g.} \emph{timeline} agregada de auditoria). & \\
\midrule
13    & 2--6~jun       & Revisão geral do sistema; validação dos critérios de aceitação. & \\
\midrule
14    & 9--13~jun      & Redacção dos Capítulos~4 e~5 do relatório final; bibliografia APA; anexos. & \\
\midrule
15    & 16--20~jun     & Reunião de preparação para defesa; ensaio de perguntas de júri. & Prep.\ defesa \\
\midrule
16    & 24~jun         & Submissão do relatório final e do código. & \textbf{Final (24~jun)} \\
\bottomrule
\caption{Calendário das 16 semanas com marcos formais sublinhados.}
\label{tab:calendario}
\end{longtable}
